\section{Conclusion}

This paper presented a comprehensive study of the stohastic Capacitated Vehicle Routing Problem with Time Windows. The formulated mathematical model includes vehicle capacity constraints, time windows, service times, and operational cost components, including travel distance, waiting time and penalty costs for time window violations.


\subsection{Key Findings}

The experimental study performed on 3,000 stochastic scenarios derived from six Solomon instances confirmed all four research hypotheses and demonstrated a clear advantage of the Tabu Search metaheuristic over the greedy baseline.

\begin{itemize}
    \item \textbf{ Superior expected performance (H1):}  
    Tabu Search achieved a significantly lower mean objective value (11286.18) compared to Greedy (13735.27).  
    Improvements are visible across the entire distribution: the median dropped from 7725.60 to 5990.30, 
    the 90th percentile from 29886.59 to 24795.62, 
    and the worst-case scenario decreased from 33634.45 to 26737.48.  
    Statistical tests (paired t-test and Wilcoxon) confirmed the significance of this improvement ($p=0.0$).
    
    \item \textbf{Variability Reduction (H2):}
    Tabu Search exhibited lower variability across scenarios (std. 8082.21 vs.\ 9393.89), 
    as well as consistently smaller quartile spreads and upper-tail percentiles.  
    ECDF plots showed that Tabu Search dominates Greedy for all Solomon families (C, R, RC), 
    confirming robustness regardless of spatial instance structure.

    \item \textbf{Time-Window Penalties (H3):} 
    Tabu Search achieved lower time-window violation penalties ($3373.77$ vs.\ $4313.23$), while simultaneously reducing travel cost ($1988.19$ vs.\ $2727.24$) and operation time ($5924.21$ vs.\ $6694.80$). 
    This confirms that the method handles temporal constraints more effectively, limiting delay propagation and yielding a significantly lower overall objective value.

    \item \textbf{Reduction of extreme worst-case outcomes (H4):}  
    Tabu Search substantially reduced the loss tail, lowering the 95th percentile from 30627.51 to 25260.80 
    and the maximum observed cost by almost 6900.  
    This behaviour is clearly visible in ECDF curves, where Tabu is steeper and reaches full cumulative mass earlier, 
    confirming that it generates far fewer high-cost scenarios.
\end{itemize}

\subsection{Contributions}

The main contributions of this work include:
\begin{enumerate}
    \item A complete MILP formulation of the CVRPTW with soft time windows that balances multiple cost components and supports practical operational constraints.
    \item A comprehensive experimental framework for evaluating routing algorithms under stochastic demand and time window perturbations, addressing realistic uncertainty in logistics operations.
    \item Empirical validation demonstrating that metaheuristic approaches (specifically Tabu Search) provide not only better average performance but also superior robustness and risk mitigation compared to constructive greedy methods.
    \item Analysis of cost component trade-offs showing how intelligent search can strategically balance competing objectives (route length, service time, and time window compliance).
\end{enumerate}

\subsection{Practical Implications}

The results have direct implications for logistics practitioners facing dynamic and uncertain operating conditions. The Tabu Search approach offers:
\begin{itemize}
    \item Significant cost savings (approximately 16\% on average) across diverse operational scenarios
    \item More predictable performance with reduced sensitivity to input perturbations
    \item Better management of worst-case scenarios, crucial for service level agreements and operational planning
\end{itemize}

\subsection{Limitations and Future Work}

While this study provides valuable insights, several directions warrant further investigation:

\begin{itemize}
    \item \textbf{Algorithm Refinement:} Exploration of hybrid metaheuristics combining Tabu Search with other operators (genetic algorithms, variable neighborhood search) may yield additional improvements.
    
    \item \textbf{Dynamic and Real-time Scenarios:} Extension to fully dynamic VRPTW where customer requests arrive during route execution, requiring real-time re-optimization.
    
    \item \textbf{Multi-objective Formulation:} Explicit treatment of competing objectives (cost, time, penalties) using Pareto-based multi-objective optimization techniques.
    
    \item \textbf{Heterogeneous Fleet:} Incorporation of vehicle heterogeneity with different capacities, costs, and operational characteristics.
    
    \item \textbf{Sustainability Considerations:} Integration of environmental objectives such as CO$_2$ emissions minimization alongside operational cost reduction.
    
    \item \textbf{Machine Learning Integration:} Use of predictive models to anticipate demand patterns and time window violations, enabling proactive route planning.
\end{itemize}

In conclusion, this work demonstrates that metaheuristic optimization approaches, particularly Tabu Search, offer substantial benefits for solving real-world vehicle routing problems under uncertainty. The combination of mathematical rigor in problem formulation, comprehensive experimental design, and practical implementation provides a solid foundation for both academic research and industrial application in logistics optimization.